% Generated draft: requires mathematical and source review.
\documentclass[11pt]{article}
\usepackage[margin=25mm]{geometry}
\usepackage{fontspec}
\setmainfont{Latin Modern Roman}
\usepackage{amsmath,amssymb,mathtools,unicode-math}
\setmathfont{Latin Modern Math}
\usepackage{enumitem,needspace,xurl,hyperref}
\setlength{\parindent}{0pt}
\setlength{\parskip}{5pt}
\setlength{\emergencystretch}{3em}
\newcommand{\N}{\mathbb{N}}
\newcommand{\Z}{\mathbb{Z}}
\newcommand{\Q}{\mathbb{Q}}
\newcommand{\R}{\mathbb{R}}
\newcommand{\C}{\mathbb{C}}
\newcommand{\F}{\mathbb{F}}
\newcommand{\A}{\mathbb{A}}
\newcommand{\PP}{\mathbb{P}}
\newcommand{\E}{\mathbb{E}}
\newcommand{\eps}{\varepsilon}
\newcommand{\dd}{\,\mathrm{d}}
\begin{document}
\section*{UnsolvedMath ID: 30001065}
\textbf{Source problem code:} catalogue entry
\subsection{Definitions and mathematical statement}
A spherical code with \(N\) points in \(\R^n\) is a set \(C\subset S^{n-1}=\{x:\|x\|=1\}\) of \(N\) distinct points. In the standard squared-chordal-distance convention, the energy for a potential \(f\) is
\[
E_f(C)=\sum_{\{x,y\}\subset C,\ x\ne y}f(\|x-y\|^2),
\]
where each unordered pair is counted once. A completely monotonic potential satisfies
\[
(-1)^j f^{(j)}(t)\geq0\quad(j\geq0)
\]
on the relevant positive squared-distance interval, conventionally \(0<t<4\) with its value at \(4\) interpreted continuously. The precise admissible potential domain and regularity convention are not given in the source.

The code \(C\) is universally optimal if, for every such potential \(f\) and every other \(N\)-point set \(D\subset S^{n-1}\),
\[
E_f(C)\leq E_f(D).
\]
The cardinality and ambient dimension are fixed in each comparison. The source asks whether the two experimentally identified balanced codes, one with \(N=40,n=10\) and one with \(N=64,n=14\), have this property. Thus their spheres are \(S^9\) and \(S^{13}\), respectively. It also asks for further universal optima beyond the known examples, in the original context of dimensions \(n\geq4\).

The actual two codes are not specified by coordinates, Gram matrices, or a construction in the extraction. Their cardinalities, dimensions, and the word ``balanced'' do not by themselves identify the intended configurations. Balanced commonly concerns equilibrium under radial pair interactions and must not be treated as an established universal-minimum property. Both the intended codes and the source's precise balanced and potential conventions must therefore be recovered before testing the conjectures. The existing list meant by ``known ones'' is also absent.
\subsection{Short English statement}
Are the experimentally found balanced \(40\)-point code in \(\R^{10}\) and \(64\)-point code in \(\R^{14}\) global energy minima for every admissible completely monotonic potential? Are there additional universal optima in the stated higher-dimensional context?
\subsection{Sources}
\begin{itemize}
\item Local UnsolvedMath registry, permanent ID \(30001065\), source code catalogue entry, ``Universal Optimality of Exceptional Spherical Codes.'' No individual proposer is supplied.
\item \emph{Discrete Geometry} (2009), Oberwolfach Reports, DOI 10.4171/owr/2008/44, pp. 2477--2556, listed as ``OWR Report 44'': \url{https://doi.org/10.4171/owr/2008/44}.
\item Supplied reference, ``Universal minima of spherical designs'': \url{https://arxiv.org/abs/2301.06623}.
\item Supplied literature triage, checked 2026-08-21: status ``open.'' It reports candidate balanced codes with universal optimality explicitly unproved and says later related-minimum work does not establish this claim. This is supplied dated curation, not independently verified; time-sensitive status and cited claims require independent verification.
\end{itemize}
\end{document}
